%!TEX program = lualatex
\documentclass[10pt, openany]{book}

\usepackage{comment}
\usepackage{silence}
\WarningFilter{latex}{Command \underbar has changed}
\WarningFilter{latex}{Command \underline has changed}

% ============== 考虑出血的版式 ===================================
\usepackage[
    paperwidth=6in,
    paperheight=9in,
    top=0.8in,
    bottom=0.8in,
    inner=0.75in,
    outer=0.6in,
    bindingoffset=0.25in,
    footskip=0.4in
]{geometry}

% ============== 语言和字体 ===================================
\usepackage[x-4]{pdfx} 
\usepackage{fontspec}
\usepackage[provide=*, chinese]{babel}
\usepackage{sectsty}

\babelfont{rm}[
    Path = ./font/,
    Renderer = HarfBuzz,
    BoldFont = {GFSDidot-Regular.ttf},
    BoldFeatures = {RawFeature={embolden=3.0}},
    ItalicFont = {timesi.ttf}
]{GFSDidot-Regular.ttf}

\babelfont{sf}[Path = ./font/, Renderer = HarfBuzz]{GFSDidot-Regular.ttf}
\babelfont{tt}[Path = ./font/, Renderer = HarfBuzz]{NotoSerifSC-Regular.ttf}

\babelfont[chinese]{rm}[
    Path = ./font/,
    Renderer = HarfBuzz,
    BoldFont = {NotoSerifSC-ExtraBold.ttf},
    ItalicFont = {StKaiTi.ttf}
]{NotoSerifSC-Regular.ttf}

\allsectionsfont{\rmfamily\bfseries}

\usepackage{amssymb}
\usepackage{enumitem}
\usepackage{titlesec}
\usepackage{quoting}
\usepackage{setspace}
\onehalfspacing
\usepackage{tikz}
\usepackage{fancyhdr}
\usepackage{amsmath, amsthm}
\usepackage{bm} 
\usepackage{libertine} 
\usepackage[libertine]{newtxmath} 
\usepackage{tikz-cd} 
\usepackage{nicematrix} 
\usepackage{adjustbox} 
\usepackage{framed} 
\usepackage{scrextend}
\changefontsizes{9.5pt} 

\NiceMatrixOptions{
  custom-line = {
    letter = : ,
    command = hdashline ,
    tikz = { dashed }
  }
}

\setlength{\parskip}{1em} 

\pagestyle{fancy}
\fancyhf{}
\fancyhead[LE,RO]{\thepage}
\fancyhead[RE]{\textbf{西 \quad 窗 \quad 烛}}
\fancyhead[LO]{\textit{The Candle by the West Window}}
\renewcommand{\headrulewidth}{0.4pt}

\begin{document}

\frontmatter
\title{\textbf{张量与矩阵}\\ \large 范畴论视角}
\author{Lao Chen}
\date{\today}
\maketitle
\tableofcontents

\mainmatter

\chapter{集合的代数}
在进入线性代数之前，我们必须审视最基础的范畴：集合范畴 $\mathbf{Set}$。

\section{势与等势性}
集合的“大小”被定义为它的势 $|A|$。若存在双射 $f: A \to B$，则 $|A| = |B|$，这等价于 $A, B$ 在 $\mathbf{Set}$ 中同构。

\section{集合的算术运算与势的保持}
集合范畴的构造完美对应了算术运算：
\begin{itemize}
    \item \textbf{加法（共积）}：$|A \sqcup B| = |A| + |B|$。
    \item \textbf{乘法（积）}：$|A \times B| = |A| \times |B|$。
    \item \textbf{指数（函数对象）}：$|A^B| = |A|^{|B|}$，表示所有 $B \to A$ 的映射。
\end{itemize}
这些构造满足分配律 $A \times (B \sqcup C) \cong (A \times B) \sqcup (A \times C)$，这为张量积对直和的分配律提供了最底层的范畴论支撑。


\chapter{张量}

\section{集合}

集合和集合间的函数映射构成\textbf{集合范畴(category of sets)}，
记为$\textbf{Set}$。
$f:X \to Y$相当于$f\in\textbf{Set}$，
也就是说函数$f$构成是范畴$\textbf{Set}$中的\textbf{箭头(arrow)}。
在范畴论中，
箭头的起止端点是范畴中的\textbf{对象(object)},
记为起点$X = \text{Dom} \ f \in \textbf{Set}$，
终点$Y = \text{Cod} \ f \in \textbf{Set}$。

若有有限集合$X = \{a,b,c\}$，$Y=\{a,e\}$，
若给出一个指标排序，则有列向量表示：

$$
X=\begin{pmatrix}x_1 \\ x_2 \\ x_3\end{pmatrix}=\begin{pmatrix}a \\ b \\ c\end{pmatrix},\ 
Y=\begin{pmatrix}y_1 \\ y_2\end{pmatrix}=\begin{pmatrix}a \\ e\end{pmatrix}
$$

这是有限集合的情况，
甚至对于无限可列集合乃至不可列集，
也可以有抽象的列的概念。
为防止混淆，
用圆括号表示集合，
用方括号表示普通意义的向量。

按顺序从$X$和$Y$中各取出一个元素，
构成\textbf{卡积(Cartesian product)}中的二元组
$(x,y) \in X \times Y$。
考虑某种称为\textbf{张量积(tensor product)}的二元（不一定交换的）运算$\otimes$，
从二元组生成一个新元素$x \otimes y\in Z$，
它属于某个目前尚不确定的集合$Z$。
最简单的例子就是对二元组不做任何处理：

$$
\begin{aligned}
\otimes: X \times Y &\to Z = X \times Y \\
(x,y) &\mapsto x \otimes y = (x,y) \\
\end{aligned}
$$

这样方式定义的$x \otimes y = (x,y)$的全体构成集合的张量积
$X \otimes Y = X \times Y$。
它保留了$X$和$Y$元素配对的所有信息，
而不施加任何额外的代数缩并，
而张量的乐趣在于在比集合更复杂的数学结构中，
可以定义各种有趣的代数缩并。
保持列向量$X$和$Y$的排序来考察$X \otimes Y$的列向量表示：

\begin{equation}
\begin{aligned}
X \otimes Y 
&=\begin{pmatrix}x_1 \\ x_2 \\ x_3\end{pmatrix} \otimes \begin{pmatrix}y_1 \\ y_2\end{pmatrix}
=\begin{pmatrix}a \\ b \\ c\end{pmatrix} \otimes \begin{pmatrix}a \\ e\end{pmatrix} \\
&=\begin{pmatrix}x_1 \\ x_2 \\ x_3\end{pmatrix} \otimes Y
=\begin{pmatrix}a \\ b \\ c\end{pmatrix} \otimes Y \\
&=\begin{pmatrix}x_1 \otimes Y \\ x_2 \otimes Y \\ x_3 \otimes Y\end{pmatrix}
=\begin{pmatrix}a \otimes Y \\ b \otimes Y \\ c \otimes Y\end{pmatrix} \\
&=\begin{pmatrix}x_1 \otimes y_1 \\x_1 \otimes y_2 \\
    x_2 \otimes y_1 \\x_2 \otimes y_2 \\
    x_3 \otimes y_1 \\x_3 \otimes y_2 \\\end{pmatrix}
=\begin{pmatrix}a \otimes a \\ a \otimes e \\
    b \otimes a \\ b \otimes e \\
    c \otimes a \\ c \otimes e \\\end{pmatrix} 
=\begin{pmatrix}(a,a) \\ (a,e) \\
    (b,a) \\ (b,e) \\
    (c,a) \\ (c,e) \\\end{pmatrix} \\
&= \{ (x_i,y_j) \ | \ 1 \leq i \leq 3,\ 1 \leq j \leq 2 \} \\
\end{aligned}
\end{equation}\label{eq:set-tensor}

显然

$$
|X \otimes Y| = |X \times Y| = |X| \times |Y|
$$

\section{形式和}

变量$x$的正整数次幂的集合为$X = \{1,x,x^2,\cdots\}$，
实域$\mathbb R$上的一元多项式是$X$按照$\mathbb R$系数的有限形式和，
例如：

$$
p(x) = 3x^3 - 2x + 5
$$

换个角度看，
$p$给$X$中的每个元素指定了一个$\mathbb R$系数：

$$
\begin{aligned}
p: X &\to \mathbb R \\
x^i &\mapsto a_i
\end{aligned}
$$

使得$p(x) = a_0 + a_1 x + a_2 x^2 + \cdots$，
且$a_i \neq 0$的系数是有限的。
从另一个角度，
可以认为$X = \{1,x,x^2,\cdots\}$是基，
通过$\mathbb R$系数生成了更大的数学结构，
按照前面讨论的集合列向量表示，
可以清晰看出这种$\mathbb R$线性组合的生成方式：

\begin{equation}
\underbrace{X=\begin{pmatrix}1 \\ x \\ x^2 \\ x^3 \\ \vdots \end{pmatrix}}_{\text{基}}
\quad \xrightarrow{\text{赋予系数}} \quad
\underbrace{p(x)=\text{sum}\begin{pmatrix}5 \\ -2 x \\ 0 \\ 3 x^3 \\ \vdots\end{pmatrix}}_{\text{形式项}}
\quad \xrightarrow{\text{提取坐标}} \quad
\underbrace{p=\begin{pmatrix}5 \\ -2 \\ 0 \\ 3 \\ \vdots\end{pmatrix}}_{\text{系数}}
\end{equation}\label{eq:polynomial-vector}

这样多项式$p$完全由一个系数列向量表示。
当这个列向量中只有有限非零元时，
对应了有限形式和。
在范畴论中，
集合$X$上取值为$\mathbb R$的函数，
即$\mathbb R$值函数全体所构成的函数集合，
常记为：

$$
\mathbb R^X = \{p: X \to \mathbb R \} 
$$

它是集合范畴中的\textbf{指数对象(exponential object)}，
读者可以自行验证无论有限无限都有$|\mathbb R^X| = |\mathbb R| ^ {|X|}$。
此外，
它还可以用列向量空间表示，
这里我们用方括号矩阵，
体现它就是我们在线性代数中讨论的那种矩阵：

$$
\mathbb R^X = \left\{ \begin{bmatrix} a_1 \\ a_2 \\ \vdots \end{bmatrix}, \quad a_i \in \mathbb R \right\}
$$

多项式集合$\mathbb R[X]$描述的是函数集合$\mathbb R^X$中，
其中满足有限非零系数的那一部分。
在这里我们看到了如何从$X = \{1,x,x^2,\cdots\}$的各次幂作为生成元，
通过不同的系数来构造多项式。
生成元按某种系数的构造，
具有列向量的表示形式，
是我们接下来要讨论的。

有了多项式这个例子的铺垫，
我们可以开始对任何集合$X\in\textbf{Set}$赋予类似的代数结构。
正如多项式需要用次数来区分生成元，
这里先给出有限或无限的指标集合$I\in\textbf{Set}$，
满足$|I| = |X|$，
使得指标$i\in I$能单满映射到集合元素$x_i\in X$。
给集合的每个元素赋予一个整数：

$$
\begin{aligned}
p: X &\to \mathbb Z \\
x_i &\mapsto p(x_i) = n_i \\
\end{aligned}
$$

宏观地看，
所有这样的赋予方式构成函数集，
用范畴论的记号记为：

$$
\mathbb Z^X = \{p: X \to \mathbb Z\} = \textbf{Set}(X,\mathbb Z) = \text{Hom}(X,\mathbb Z) 
$$

类似前面在多项式集合的讨论，
这里指数形式暗示它是一个\textbf{指数对象(exponential object)}，
$\textbf{Set}(-,-)$或$\text{Hom}(-,-)$则给出了一个集合范畴中的\textbf{态射集(hom-set)}也就是箭头集。
而以上的指定方式是这个态射集中的成员函数$p\in \mathbb Z^X$，
这样对应到集合$X$按照整系数的\textbf{形式和(formal sum)}：

$$
\sum_{i\in I} n_i x_i
$$

本质上其中的项是$p$构造的输入输出的二元组$(x_i,n_i)$，
从而构成集合：

$$
p_X = \{ (x_i,n_i) \}_{i\in I}
$$

从另一个角度，
可以认为$X$是基，
通过$\mathbb Z$系数生成了更大的数学结构，
类似(\ref{eq:polynomial-vector})：

\begin{equation}
\underbrace{X=\begin{pmatrix}x_1 \\ x_2 \\ x_3\end{pmatrix}}_{\text{基}}
\quad \xrightarrow{\text{赋予系数}} \quad
\underbrace{p_X=\text{sum}\begin{pmatrix}n_1 x_1 \\ n_2 x_2 \\ n_3 x_3\end{pmatrix}}_{\text{形式项}}
\quad \xrightarrow{\text{提取坐标}} \quad
\underbrace{p=\begin{pmatrix}n_1 \\ n_2 \\ n_3\end{pmatrix}}_{\text{系数}}
\end{equation}\label{eq:formal-sum-vector}

同样可以讨论用方括号矩阵表示的列向量空间：

$$
\mathbb Z^X = \left\{ \begin{bmatrix} n_1 \\ n_2 \\ \vdots \end{bmatrix}, \quad n_i \in \mathbb Z \right\}
$$

注意到向量空间的维度和集合势的关系：

\begin{equation}
\dim \mathbb Z^X = |X|
\end{equation}\label{eq:dim-card-Z}

类似多项式的处理，
通常当$n_i=0$时，
$n_i x_i$这一项要消掉，
从形式求和公式中排除。
当我们讨论\textbf{有限形式和(finite formal sum)}时，
要求上式子中$n_i \neq 0$的情况为有限个，
排除了零系数项后，
得到有限个非零系数的形式和。
相应的，
类似多项式集合的记号，
有限形式和用指数对象的方式记为$\mathbb Z[X]$，
它是函数集$\mathbb Z^X$中有限非零系数的那部分，
是它的子空间，
因为它对加法运算封闭，
用列向量表示：

$$
\begin{aligned}
\mathbb Z[X] \times \mathbb Z[X] &\to \mathbb Z[X] \\    
\left(
    \begin{bmatrix} m_1 \\ m_2 \\ \vdots \end{bmatrix},
    \begin{bmatrix} n_1 \\ n_2 \\ \vdots \end{bmatrix}
\right) 
&\mapsto
    \begin{bmatrix} m_1 \\ m_2 \\ \vdots \end{bmatrix} +
    \begin{bmatrix} n_1 \\ n_2 \\ \vdots \end{bmatrix}
    = \begin{bmatrix} m_1 + n_1 \\ m_2 + n_2 \\ \vdots \end{bmatrix}
\end{aligned}
$$

有限非零个$m_i$和$n_j$的相加，
得到有限非零个$m_i + n_j$。
这使得我们完全可以用线性代数的角度讨论它。

\section{自由函子}

前面讨论了集合$X$的整系数有限形式和全体$\mathbb Z[X]$，
显然可以定义其它的系数，
如环或者域，
关键是这个系数集合$\Lambda$中要有$0$的概念，
才能消去形式项，
满足有限性的要求，
形成$\Lambda$系数的有限形式和全体$\Lambda[X]$。
这里的$0$的概念暗示了Abel群的单位元，
和相应的具有交换律的加性代数结构。

以$\mathbb Z[X]$为例，
$\mathbb Z = (\mathbb Z,+,0)\in\textbf{Ab}$是Abel群，
$\textbf{Ab}$是Abel群范畴，
$\mathbb Z$的Abel群结构诱导出$\mathbb Z[X]$上的Abel群的结构。
用函子描述：

$$
\begin{aligned}
F: \textbf{Set} &\to \textbf{Ab} \\
X &\mapsto FX = \mathbb Z[X] \\
Y &\mapsto FY = \mathbb Z[Y] \\
\end{aligned}
$$

注意到我们只讨论过一种二元集合运算，
就是集合的张量积$X \otimes Y = X \times Y$，
我们自然关注$F(X \otimes Y) = \mathbb Z[X \otimes Y]$能得到什么，
它和$FX = \mathbb Z[X]$以及$FY = \mathbb Z[Y]$有什么关系。
根据(\ref{eq:dim-card-Z})有：

\begin{equation}
\dim \mathbb Z^{X \otimes Y} = \dim \mathbb Z[X \otimes Y] = |X \otimes Y| = |X| \times |Y|
= \dim \mathbb Z^X \times \dim \mathbb Z^Y
\end{equation}\label{eq:dim-card}

再考虑到势：

$$
% https://tikzcd.yichuanshen.de/#N4Igdg9gJgpgziAXAbVABwnAlgFyxMJZABgBpiBdUkANwEMAbAVxiRAA0ACAHW4jwC28TgE1OAXk5deg4SJABfUuky58hFABZyVWoxZteAujgAWAI3OcAWsml9ZcURUXKQGbHgJFtm3fWZWRBBeKCwBHm5jM0sbO0j+cLkKCU4AH3tEoScRNNSMvJkkpzTcxV0YKABzeCJQADMAJwgBJDIQHAgkAEZqAINggDFXBubWxF6OrsQAJj79IJDuMNalUZa26k6kOb1AtjTeAGMofjSQagY6cxgGAAVVLw0QRqwq0xxyhSA
\begin{tikzcd}
X \otimes Y = X \times Y \arrow[rrrr, "F"] \arrow[rrrrdddd, "|\cdot|"'] &  &  &  & {\mathbb Z[X \otimes Y]} \arrow[dddd, "\dim"]                  \\
                                                                        &  &  &  &                                                                \\
                                                                        &  &  &  &                                                                \\
                                                                        &  &  &  &                                                                \\
                                                                        &  &  &  & {\dim \mathbb Z[X \otimes Y] = |X \otimes Y| = |X| \times |Y|}
\end{tikzcd}
$$

根据(\ref{eq:set-tensor})：

$$
F(X \otimes Y)
=F\begin{pmatrix}a \otimes a \\ a \otimes e \\ b \otimes a \\ b \otimes e \\ c \otimes a \\ c \otimes e\end{pmatrix}
$$

类似(\ref{eq:formal-sum-vector})：

\begin{equation}
\underbrace{X \otimes Y=\begin{pmatrix}a \otimes a \\ a \otimes e \\ b \otimes a \\ b \otimes e \\ c \otimes a \\ c \otimes e\end{pmatrix}}_{\text{基}}
\quad \xrightarrow{\text{赋予系数}} \quad
\underbrace{p_{X \otimes Y}=\text{sum}\begin{pmatrix}n_{a \otimes a}(a \otimes a) \\ n_{a \otimes e}(a \otimes e) \\ n_{b \otimes a}(b \otimes a) \\ n_{b \otimes e}(b \otimes e) \\ n_{c \otimes a}(c \otimes a) \\ n_{c \otimes e}(c \otimes e)\end{pmatrix}}_{\text{形式项}}
\quad \xrightarrow{\text{提取坐标}} \quad
\underbrace{p=\begin{pmatrix}n_{a \otimes a} \\ n_{a \otimes e} \\ n_{b \otimes a} \\ n_{b \otimes e} \\ n_{c \otimes a} \\ n_{c \otimes e}\end{pmatrix}}_{\text{系数}}
\end{equation}

系数$p$的全体构成了维度为$|X| \times |Y|$的向量空间(\ref{eq:dim-card})，
即张量空间。
注意到函子对张量积的保持：

$$
\begin{aligned}
F: \textbf{Set} &\to \textbf{Ab} \\
X &\mapsto FX = \mathbb Z[X] \\
Y &\mapsto FY = \mathbb Z[Y] \\
X \otimes Y &\mapsto F(X \otimes Y) = \mathbb Z[X] \otimes \mathbb Z[Y] \\
\end{aligned}
$$

\section{商}

前面讨论过一个变元的多项式，
现在考虑$\mathbb R[x]$和$\mathbb R[y]$中的两对多项式：

$$
\begin{cases}
p_1(x) = 3x + 1 \\
q_1(y) = - 2y
\end{cases} \quad
\begin{cases}
p_2(x) = -6x - 2 \\
q_2(y) = y
\end{cases}
$$

从集合的张量积角度看，
$(p_1,q_1) = p_1 \otimes q_1 \in \mathbb R[x] \otimes \mathbb R[y]$，
$(p_2,q_2) = p_1 \otimes q_1 \in \mathbb R[x] \otimes \mathbb R[y]$，
且$(p_1,q_1) \neq (p_2,q_2)$。
然而在二元多项式集$\mathbb R[x,y]$中：

$$
p_1(x) \cdot q_1(y) = p_2(x) \cdot q_2(y) = - 6xy - 2y 
$$

这意味着$(p_1,q_1)$和$(p_2,q_2)$借助多项式的乘法运算建立了某种联系，
而这种联系体现了某种冗余信息。
从代数结构精细化的角度看，
我们要的不止是集合范畴的张量积$\mathbb R[x] \otimes_\textbf{Set} \mathbb R[y]$，
而是需要某种更精细的范畴中的张量积，
例如Abel群范畴中的张量积$\mathbb R[x] \otimes_\textbf{Ab} \mathbb R[y]$。
我们需要某种类似于上述$p_1(x) \cdot q_1(y) = p_2(x) \cdot q_2(y) $的性质来构造等价关系
$(p_1,q_1) \sim (p_2,q_2)$，
产生我们需要的张量积：

$$
\mathbb R[x] \otimes_\textbf{Ab} \mathbb R[y]
= \mathbb R[x] \otimes \mathbb R[y]/\sim
$$

具体而言：

$$
(p_1,q_1) \sim (p_2,q_2) \ \Longleftrightarrow \ 
p_1 \cdot q_1 = p_2 \cdot q_2 \ \Longleftrightarrow \ 
p_1 \cdot q_1 - p_2 \cdot q_2 = 0
$$


\section{自由群}

集合$X=\{a,b,c\}$作为字母表可以生成字符串，
所有由$X$生成的字符串构成的集合记$W$。
$W$上用字符串拼接可以定义封闭的加法$[aab] + [cb] = [aabcb]$，
加法有唯一的单位元即空串$[]$，
这样使得$(W,+,[])$构成非交换的含幺加法系统。

为了得到具有对称性的群结构，
将$X$的所有元素都赋予一个形式化的逆元，
字母表扩展为$\{a,a^{-1},b,b^{-1},c,c^{-1}\}$，
在此基础上得到扩展的加法系统$(FX,+,[])$。
约定$[xx^{-1}] = [x^{-1}x] = []$，
使得字符串中相邻互逆的字母被约化，
且所有的约化字符串都有唯一的逆，
这样使得$(FX,+,[])$构成群，
称为$X$生成的\textbf{自由群(free group)}。




\section{整数}

整数集$\mathbb Z$具有Abel群和交换环的结构。
给定整数$n \in \mathbb Z$有正规子群$n\mathbb Z = \{nk \ | \ k\in \mathbb Z\}$，
若将其记为$(n)$则强调它还是整数环中的理想。
同余关系的等价形式有 $a - b \in n\mathbb Z \Leftrightarrow a \sim_n b \Leftrightarrow a \equiv b \pmod n$，
得到商群：
$\mathbb Z_n = \mathbb Z/n\mathbb Z = \mathbb Z/\sim_n = \{ i + n\mathbb Z\} = \{[0]_n,[1]_n,\cdots,[n-1]_n\}$，
为简洁也可以直接记作$\mathbb Z_n = \{0,1,\cdots,n-1\}$，
这是有限集$|\mathbb Z_n| = n$。
若将其记为$\mathbb Z_n = \mathbb Z/(n)$则强调它还是整数环中的商环，
有对应的乘法$[a] \cdot [b] = [a \cdot b]$。
用前面的列向量有序表示，
构造集合的张量积：

$$
\begin{aligned}
\mathbb Z_m \otimes \mathbb Z_n 
&=\begin{pmatrix}[i]_m\end{pmatrix}
\otimes \begin{pmatrix}[j]_n\end{pmatrix} \\
&=\begin{pmatrix}[i]_m \otimes [j]_n\end{pmatrix}
=\begin{pmatrix}(i + m\mathbb Z) \otimes (j + n\mathbb Z)\end{pmatrix}
\end{aligned}
$$

在集合范畴中，张量积就是二元组$x \otimes y = (x,y)$。
现在在Abel群范畴，
我们希望张量积$(i + m\mathbb Z) \otimes (j + n\mathbb Z)$保持一些代数结构，
而不仅仅是二元组$(i + m\mathbb Z,j + n\mathbb Z)$的罗列。
自然的想法是，
既然$\mathbb Z_m$和$\mathbb Z_n$都是同余类构成的商群，
我们希望张量积空间$\mathbb Z_m \otimes \mathbb Z_n$也具有类似的形式，
这样可以在代数上封闭，
即存在$o\in\mathbb Z^+$使得
$\mathbb Z_m \otimes \mathbb Z_n \simeq \mathbb Z_o$，
于是有$(i + m\mathbb Z) \otimes (j + n\mathbb Z)$和$k + o\mathbb Z$的对应。

先考虑$n=1$时$\mathbb Z_1 = \mathbb Z/1\mathbb Z = \{0\}$是平凡群，
于是

$$
\mathbb Z_m \otimes \mathbb Z_1 = \mathbb Z_m \otimes \{0\}
= \begin{pmatrix}(i + m\mathbb Z) \otimes 0 \end{pmatrix}
\simeq \begin{pmatrix}i + m\mathbb Z\end{pmatrix}
= \mathbb Z_m
$$

进而考虑

$$
\mathbb Z_m \otimes \mathbb Z_{mn}
= \begin{pmatrix}(i + m\mathbb Z) \otimes (i + mn\mathbb Z) \end{pmatrix}
\simeq \begin{pmatrix}i + m\mathbb Z\end{pmatrix}
= \mathbb Z_m
$$


不难理解：

$$
\mathbb Z_m \otimes \mathbb Z_n \simeq \mathbb Z_{\text{gcd}(m,n)}​
$$

\begin{center}
\begin{tikzcd}[column sep=large]
(A \oplus B) \times C \arrow[r, "\phi"] \arrow[d, "f"] & (A \oplus B) \otimes C \arrow[dl, "\exists ! \bar{f}", dashed] \\
(A \otimes C) \oplus (B \otimes C) & 
\end{tikzcd}
\end{center}

\chapter{张量函子与伴随}

\section{张量-同态伴随 (Tensor-Hom Adjunction)}
张量积的核心逻辑在于它与内合同态的伴随关系：
\begin{equation}
(V \otimes -) \dashv \mathrm{Hom}(V, -)
\end{equation}
这意味着存在自然同构：$\mathrm{Hom}(V \otimes W, U) \cong \mathrm{Hom}(V, \mathrm{Hom}(W, U))$。

\section{普遍性质}
张量积将多线性映射 $B: V \times W \to U$ 唯一地线性化为 $\tilde{B}: V \otimes W \to U$：
\begin{center}
\begin{tikzcd}[column sep=large, row sep=large]
V \times W \arrow[r, "\otimes"] \arrow[rd, "B"'] & V \otimes W \arrow[d, "\exists! \tilde{B}", dashed] \\
& U
\end{tikzcd}
\end{center}

\end{document}